\documentclass[a4paper,10pt]{article} \usepackage{anysize}
\marginsize{2cm}{2cm}{1cm}{1cm}
%\textwidth 6.0in \textheight = 664pt
\usepackage{xltxtra}
\usepackage{xunicode}
\usepackage{graphicx}
\usepackage{color}
\usepackage[usenames,dvipsnames]{xcolor}
\usepackage{xgreek}
\usepackage{fancyvrb}
\usepackage{minted}
\usepackage{listings}
\usepackage{enumitem} \usepackage{framed} \usepackage{relsize}
\usepackage{float} \setmainfont[Mapping=TeX-text]{FreeSerif}
\begin{document}

\renewcommand{\theenumi}{\roman{enumi}}
\begin{titlepage}
\begin{center}
\begin{figure}[t] 
     \includegraphics[scale=0.7]{title/ntua_logo}
\end{figure}
\begin{LARGE}\textbf{ΕΘΝΙΚΟ ΜΕΤΣΟΒΙΟ ΠΟΛΥΤΕΧΝΕΙΟ\\}\end{LARGE}
\vspace{2cm}
\begin{Large}
ΣΧΟΛΗ ΗΜ\&ΜΥ\\
ΤΟΜΕΑΣ ΤΕΧΝΟΛΟΓΙΑΣ ΕΠΙΚΟΙΝΩΝΙΩΝ, ΗΛΕΚΤΡΟΝΙΚΗΣ ΚΑΙ ΣΥΣΤΗΜΑΤΩΝ ΠΛΗΡΟΦΟΡΙΚΗΣ\\
ΕΡΓΑΣΤΗΡΙΟ ΨΗΦΙΑΚΩΝ ΣΥΣΤΗΜΑΤΩΝ\\
1\textsuperscript{η} ΑΣΚΗΣΗ\\
Ακ. έτος 2010-2011\\
\end{Large}
\vspace{5cm}
\Large Τμήμα Γ, Ομάδα 2\textsuperscript{η}\\
\vspace{1cm}
\begin{tabular}{l r}
\Large{Γερακάρης Βασίλης}&
\large{Α.Μ.: 03108137}\\
\Large{Λύρας Γρηγόρης}&
\large{Α.Μ.: 03109687}\\
\end{tabular}\\
\vspace{5cm}

\vfill
\large\today\\
\end{center}
\end{titlepage}


\section*{Ερωτήσεις}
\setcounter{section}{3}
\subsection{Άσκηση 1.1}
\subsubsection{}
Στην περίπτωση που έρθει ένα σήμα κατά την εκτέλεση της ρουτίνας ενός signal
handler, αυτό αγνοείται προσωρινά, αφού γίνονται masked όταν ένας από τους 2
τρέχει.
\subsubsection{}
Τα SIGCHLD συνήθως προέρχονται από running παιδιά που δέχονται SIGSTOP από τον
scheduler όταν τελειώσει το κβάντο χρόνου τους (οπότε φέρνουμε για εκτέλεση
την επόμενη διεργασία) ή SIGCONT (όταν μια ready διεργασία επιλεγεί για
εκτέλεση, που το αγνοούμε). Επίσης έρχεται 1 σήμα SIGCHILD (κατα τη δημιουργία
μιας διεργασίας, όπου κάνει raise SIGSTOP στον εαυτό της) και 1 όταν κάνει
έξοδο. Με τη χρήση μιας global μεταβλητής current\_proc, όπου διατηρούμε το pid
της διεργασίας που εκτελείται κάθε στιγμή, μπορούμε να
\subsubsection{}
Μια διεργασία έχει τη δυνατότητα να πιάσει το σήμα SIGSTΟP. Αν συμβεί αυτό,
ένας scheduler που χρησιμοποιεί μόνο SIGALRM και δεν ελέγχει αν η διεργασία
που προσπαθεί να σταματήσει όντως σταμάτησε, είναι δυνατόν να φτάσει σε
κατάσταση όπου περισσότερες από μία διεργασίες προσπαθούν να τρέξουν. Επίσης
το SIGCHILD χρησιμοποιείται για να ειδοποιηθεί ο scheduler σε περίπτωση που
κάποιο παιδί του πεθάνει (από κάποιο σήμα είτε από κανονική έξοδο).

\subsection{Άσκηση 1.2}
\subsubsection{}
Εμφανίζεται πάντα ο φλοιός ως current\_proc, αφού προκειμένου να εκτελέσει τις
εντολές που του δώθηκαν (μαζί και την 'p'), πρέπει να είναι running. Δε
γίνεται να αποφευχθεί αυτό, αφού οι εντολές του πρέπει αυστηρά να εκτελούνται
μέσα στο κβαντο χρόνου που του αντιστοιχεί.
\subsubsection{}
Αν δε γινόταν απενεργοποιήση των σημάτων πριν την υλοποίηση των συναρτήσεων
του φλοιού θα ήταν πιθανό να έρθει ένα σήμα που θα μεταβάλει την κατάσταση της
ουράς εκτέλεσης την ώρα που ο φλοιός (για παράδειγμα) διέγραφε μια διεργασία.
Αυτό το race condition μπορούσε να προκαλέσει σφάλματα όπως διαγραφή λάθος
διεργασίας, ή (χειρότερα) κατάρρευση της ουράς με λάθος ανάθεση των pointers
στις επόμενες διεργασίες.

\subsection{Άσκηση 1.3}
\subsubsection{}
Θα ήταν δυνατόν να λιμοκτονήσει μια low\_priority διεργασία αν συνεχώς
δημιουργούσαμε και θέταμε ως high priority διεργασίες. Αφού η λίστα των υψηλής
προτεραιότητας διεργασιών θα είχε πάντα στοιχεία, ο χρονοδρομολογητής δε θα
επέλεγε ποτε τις διεργασίες με χαμηλή προτεραιότητα.

\section*{Προαιρετικές Ερωτήσεις}
\setcounter{section}{4}
\setcounter{subsection}{0}
\subsection{}
Η επικοινωνία γίνεται με τη χρήση σωληνώσεων (pipes). Ο πατέρας (scheduler)
περιμένει στο άκρο ανάγνωσης του pipe του παιδιού για να λάβει κάποια αίτηση
(request) όπως αυτή ορίζεται στο request.h. Όταν τη λάβει, απενεργοποιεί τα σήματα,
και την εξυπηρετεί. Στη συνέχεια ενεργοποιεί και πάλι τα σήματα και γράφει στο άκρο
εγγραφής του δικού του pipe. Στην άλλη μεριά το shell περιμένει στο άκρο
ανάγνωσης του pipe του πατέρα για να λάβει την απόκριση του scheduler.
Αφού τη λάβει συνεχίζει κανονικά τη λειτουργία του.

\subsection{}
Θα ήταν δυνατόν να διατηρείται στη λίστα με τις low\_priority διεργασίες ένας
counter που θα δείχνει πόσα κβάντα χρόνου έχουν περάσει από την δημιουργία
τους, ενώ υπάρχουν high\_priority διεργασίες. Ο counter αυτός θα αυξανόταν μέσα
στον SIGALRM handler (δίνοντας βέβαια αρκετή καθυστέρηση στο πρόγραμμα). Όταν
μία διεργασία έφτανε ένα προκαθορισμένο αριθμό στον counter, θα ανέβαινε σε
high\_priority ώστε να αρχίζει να εκτελείται.
\subsection{}
\begin{enumerate}
    \item
    \item
    \item
    \item
\end{enumerate}


\vspace{1cm}
\def\thesubsection {Άσκηση \arabic{section}.\arabic{subsection}}
\section*{Πηγαίος κώδικας ασκήσεων}
\emph{Κοινή βιβλιοθήκη για διαχείρηση διπλά συνδεδεμένης, κυκλικής λίστας}
\inputminted[linenos,fontsize=\footnotesize,frame=leftline]{c}{files/queue.h}
\inputminted[linenos,fontsize=\footnotesize,frame=leftline]{c}{files/queue.c}
\setcounter{section}{1}
\setcounter{subsection}{0}
\subsection{}
\inputminted[linenos,fontsize=\footnotesize,frame=leftline]{c}{files/scheduler_1.1.c}
\subsection{}
\inputminted[linenos,fontsize=\footnotesize,frame=leftline]{c}{files/scheduler-shell_1.2.c}
\subsection{}
\inputminted[linenos,fontsize=\footnotesize,frame=leftline]{c}{files/scheduler-shell_1.3.c}

\section*{Παραδείγματα εκτέλεσης των προγραμμάτων}
\def\thesubsection {Πρόγραμμα \arabic{section}.\arabic{subsection}:}
\setcounter{subsection}{0}
\subsection{Scheduler}
\lstset{emph={PID, NMSG, delay, id},emphstyle=\color{blue}}
\lstset{emph={SIREN},emphstyle=\color{yellow}}
\lstset{emph={DIE, DEAD, STOP},emphstyle=\color{red}}
\lstset{emph={NEXT, NEEEEEEEEEEXT},emphstyle=\color{ForestGreen}}
\lstinputlisting[numbers=left,numberstyle=\tiny,basicstyle=\footnotesize\ttfamily,breaklines=true]{files/scheduler_1.1.out}
%\inputminted[linenos,fontsize=\footnotesize,frame=leftline]{bash}{files/scheduler_1.1.out}
\subsection{Scheduler with shell}
\lstinputlisting[numbers=left,numberstyle=\tiny,basicstyle=\footnotesize\ttfamily,breaklines=true]{files/scheduler-shell_1.2.out}
%\inputminted[linenos,fontsize=\footnotesize,frame=leftline]{bash}{files/scheduler-shell_1.2.out}
\subsection{Scheduler with shell \& priorities}
\lstinputlisting[numbers=left,numberstyle=\tiny,basicstyle=\footnotesize\ttfamily,breaklines=true]{files/scheduler-shell_1.3.out}
%\inputminted[linenos,fontsize=\footnotesize,frame=leftline]{bash}{files/scheduler-shell_1.3.out}

\end{document}
